%!TeX root=MemoriaTFG.tex

\chapter{Discusión}
\label{cap:discusion}

\section{Significado de los resultados}
\label{sec:significado}

La lectura agregada de los experimentos no apunta a la existencia de
un único modelo vencedor, sino a que el rendimiento dermatológico está
condicionado por cuatro dimensiones principales: el nivel ontológico al
que se evalúa, el régimen de datos etiquetados disponible, la modalidad
de entrada y la estrategia de adaptación empleada. Los resultados
sugieren que, una vez alcanzado un determinado nivel de capacidad
representacional en los codificadores visuales modernos, las diferencias
de rendimiento dependen cada vez más de factores relacionados con los
datos, la formulación de la tarea y la integración multimodal. Entre
ellos destacan la taxonomía utilizada para organizar las clases, el
volumen y la calidad de las etiquetas, el alineamiento texto-imagen de
las ramas contrastivas, la formulación de los \emph{prompts}, la equidad
entre subgrupos y la propia estrategia de despliegue. Las subsecciones
siguientes interpretan estos factores a partir de la evidencia empírica
del capítulo~\ref{cap:resultados}, mientras que la
sección~\ref{sec:implicaciones} deriva sus implicaciones operativas para
el diseño de sistemas dermatológicos basados en modelos fundacionales.

\subsection{Ausencia de un método óptimo único: especialización por tarea}
\label{sec:disc-especializacion}

La síntesis transversal de los experimentos sobre DermapixelAI~1.0
(tabla~\ref{tab:dermapixel-extra-sintesis}) muestra que no existe un
método universalmente superior para todas las tareas evaluadas. Los
distintos modelos alcanzan su mejor rendimiento en escenarios diferentes
según el nivel ontológico, la modalidad de entrada y el criterio de
evaluación considerado. Así, los codificadores especializados destacan en
transferencia visual y clasificación supervisada, las arquitecturas
contrastivas muestran ventajas en recuperación multimodal y búsqueda
semántica, los métodos basados en recuperación resultan competitivos en
escenarios de cola larga diagnóstica y los \emph{ensembles} aportan
sensibilidad adicional en tareas de seguridad clínica como el cribado de
melanoma.

Esta diversidad de resultados sugiere que el problema no puede
formularse en términos de un único modelo vencedor. La elección óptima
depende del objetivo perseguido: clasificación general, subcategorización
clínica, recuperación de casos similares, segmentación, razonamiento
multimodal o maximización de sensibilidad diagnóstica. En consecuencia,
la pregunta relevante deja de ser \emph{qué modelo es el mejor} en
términos absolutos y pasa a ser \emph{qué combinación de codificador,
estrategia de adaptación y mecanismo de inferencia resulta más adecuada
para una tarea concreta}.

Este hallazgo conecta con una idea recurrente a lo largo del
capítulo~\ref{cap:resultados}: a medida que los modelos fundacionales
alcanzan niveles elevados de capacidad representacional, las diferencias
de rendimiento dependen cada vez más de cómo se formula el problema, de
la estructura de los datos y del objetivo clínico perseguido que de la
arquitectura aislada del modelo. Las subsecciones siguientes desarrollan
esta conclusión desde tres perspectivas complementarias: el nivel
ontológico de clasificación, la estrategia de adaptación y el paradigma
de inferencia empleado.

\subsubsection{Especialización por nivel ontológico}
\label{sec:disc-nivel-ontologico}

El peso del preentrenamiento de dominio queda establecido por el
sondeo lineal (sección~\ref{sec:rep-panderm},
sección~\ref{sec:benchmark}): PanDerm Large supera a PanDerm Base con
una mejora media de $+9{,}0$\,pp en BAcc sobre los diez datasets
externos contrastados (tabla~\ref{tab:lp-panderm}). El detalle por
dataset matiza el origen del beneficio: es modesto sobre dermatoscopia
estandarizada ---donde una arquitectura convolucional supervisada
sobre ImageNet-21K llega a empatar la AUROC--- y crece sobre
fotografía clínica con dispositivos heterogéneos. Los resultados
sugieren que la ventaja del preentrenamiento de dominio aumenta con la
distancia entre el dataset y los datasets web generalistas, lo que la
hace especialmente relevante en condiciones clínicas no
estandarizadas, como la atención primaria o la teledermatología.
SigLIP-Large SO400M confirma esta lectura desde el extremo opuesto:
lidera sobre HAM10000 (tabla~\ref{tab:benchmark}) por la
sobrerrepresentación de dermatoscopia estandarizada en su
preentrenamiento web, pero esa superioridad se desvanece sobre el
material clínico heterogéneo.

Esta ventaja, sin embargo, no es uniforme entre niveles ontológicos.
Sobre la cardinalidad clínica intermedia L2, DermLIP v2 se sitúa como
la mejor cabeza individual, con una ventaja de $+6{,}7$\,pp de AUROC
sobre PanDerm Large (sección~\ref{sec:dermapixel-ensemble}), y lidera
también la AUROC L2 en régimen \emph{zero-shot} sin entrenar cabeza
alguna (sección~\ref{sec:dermapixel-zs}). Ambos resultados sugieren
que el preentrenamiento contrastivo texto-imagen de
Derm1M~\cite{derm1m} captura la granularidad de subcategoría clínica
con mayor fidelidad que el preentrenamiento auto-supervisado puro de
la familia PanDerm. La jerarquía PanDerm Large $>$ DermLIP v2, válida
en la mayoría de datasets, se invierte por tanto en L2: el codificador
óptimo depende del nivel ontológico objetivo.

\subsubsection{Adaptación eficiente frente a ajuste denso}
\label{sec:disc-lora-vs-ft}

El ajuste fino parcial sobre DermapixelAI~1.0
(sección~\ref{sec:dermapixel-ft-l1}), con un volumen de entrenamiento
reducido ($N_{\mathrm{train}} = 874$), mejora sustancialmente la BAcc
sobre el sondeo lineal congelado, pero no supera al \emph{ensemble}
lineal de codificadores congelados, que lo aventaja en BAcc y AUROC
(tabla~\ref{tab:dermapixel-ensemble}). El patrón es coherente con la
literatura sobre régimen de bajo $N$: combinar linealmente
codificadores fundacionales congelados aprovecha la diversidad de sus
representaciones sin añadir los parámetros que el ajuste denso debe
estimar sobre un dataset reducido.

La cabeza supervisada L2 de Dermapixel R0
(sección~\ref{sec:dermapixel-spanderm-v0}) confirma esta lectura y
aporta la cifra-ancla del argumento: la adaptación LoRA sobre los dos
últimos bloques de PanDerm Large ---$524\,K$ parámetros entrenables,
una razón aproximada de $48\times$ frente a los $25{,}2$\,M del ajuste
denso parcial--- alcanza la mejor BAcc L2 documentada sobre el
dataset, y su AUROC L2 empata estadísticamente con la de DermLIP v2
nativo (tabla~\ref{tab:dermapixel-ft-multilevel}). El ajuste denso,
pese a su mayor parametrización, no supera a la adaptación de bajo
rango sobre el mismo encoder; sobre L3, además, ningún protocolo de
ajuste denso rebasa el techo que impone la cola larga estructural. La
implicación de diseño es directa y motiva la arquitectura de
Dermapixel R0 (sección~\ref{sec:linea-spanderm}): en régimen de bajo
$N$, la adaptación de bajo rango y la combinación de codificadores con
cabezas ligeras son preferibles al ajuste fino denso clásico.

Más allá de sus métricas concretas, Dermapixel R0 constituye una prueba
de viabilidad de la adaptación eficiente de modelos fundacionales en
entornos hospitalarios. Los resultados muestran que es posible construir
un modelo dermatológico adaptado al castellano mediante técnicas PEFT
sobre pesos previamente entrenados, sin necesidad de reproducir la
escala de datos ni de cómputo del preentrenamiento original. Mientras
que el entrenamiento fundacional sigue estando reservado a grandes
consorcios de investigación, la evidencia obtenida indica que
instituciones sanitarias y grupos académicos pueden desarrollar
adaptaciones especializadas con recursos sustancialmente más modestos.
En este sentido, Dermapixel R0 debe interpretarse menos como un sistema
de máximo rendimiento y más como una demostración de que la
especialización dermatológica de bajo coste resulta técnicamente
factible.

\subsubsection{Retrieval frente a clasificación paramétrica}
\label{sec:disc-retrieval}

Una tercera fuente de especialización emerge sobre la cola larga, en
el paradigma de clasificación. Dos observaciones la sostienen.
Primero, CLIP-L-336, un codificador generalista sin material
biomédico en su preentrenamiento, encabeza la AUROC bajo sondeo lineal
sobre L2 y L3 (sección~\ref{sec:dermapixel-extra-encoders}), lo que
matiza la hipótesis de especialización de dominio cuando la métrica es
el ranking probabilístico. Este resultado no contradice la ventaja
global del preentrenamiento dermatológico, sino que pone de manifiesto
que la AUROC y la BAcc capturan propiedades distintas del espacio de
representación: la primera mide la ordenación probabilística y la
segunda el acierto equilibrado por clase. Segundo, sustituir el clasificador
paramétrico por \emph{retrieval} $k$-NN con índice FAISS sobre los
mismos \emph{embeddings} de PanDerm Large
(sección~\ref{sec:dermapixel-extra-faiss}) supera al sondeo lineal en
el régimen de cola larga severa L3, donde el escaso número de muestras
por clase penaliza a los clasificadores paramétricos. La recuperación
por vecindad, además, es operativamente atractiva: no exige reentrenar
una cabeza al incorporar clases nuevas y permite mostrar los casos
similares que sustentan la predicción. Ambas observaciones cierran la
lectura: la elección óptima debe condicionarse al nivel ontológico y
al criterio operativo predominante, no a un codificador o clasificador
presupuesto como universal. Esta lectura sobre L3 debe entenderse, no
obstante, como indicio y no como ordenación cerrada: la cola larga
extrema y el tamaño reducido del conjunto de test en ese nivel
(sección~\ref{sec:lim-dermapixel}) limitan la fuerza de cualquier
ranking sobre L3.

\subsection{DermapixelAI 1.0 como contribución metodológica}
\label{sec:disc-dermapixel-contrib}

Buena parte de la lectura anterior se apoya en DermapixelAI~1.0, cuyo
papel en este trabajo conviene precisar: no es un dataset auxiliar de
validación, sino una contribución metodológica en sí misma. Se trata de
un dataset clínico construido en castellano, independiente del dataset de
preentrenamiento Derm1M y auditado por hash MD5 frente a él
(sección~\ref{sec:auditoria-resultados}), lo que permite evaluar
transferencia, \emph{zero-shot}, adaptación LoRA y \emph{retrieval} sin la
incertidumbre de solapamiento que afecta a otros conjuntos. Su
organización conforme a la ontología jerárquica L1/L2/L3 habilita la
comparación de modelos a varios niveles de granularidad sobre un mismo
material, y su disponibilidad documentada lo sitúa como punto de partida
para evaluar modelos fundacionales sobre dermatología clínica en español.

Más aún, DermapixelAI~1.0 no se limita a proporcionar un nuevo conjunto
de imágenes dermatológicas, sino que habilita preguntas experimentales
que los \emph{benchmarks} clásicos no permiten abordar. La estructura
jerárquica L1/L2/L3 posibilita estudiar el comportamiento de los modelos
a diferentes niveles de granularidad diagnóstica; la disponibilidad
simultánea de imagen y texto clínico en castellano permite evaluar
recuperación multimodal y alineamiento \emph{cross-modal}; y la auditoría
explícita frente a Derm1M garantiza que los experimentos de
transferencia, \emph{zero-shot} y adaptación se realizan en un escenario
libre de contaminación conocida. En este sentido, DermapixelAI~1.0 actúa
simultáneamente como \emph{benchmark}, banco de pruebas para modelos
fundacionales dermatológicos y demostrador metodológico para futuros
estudios en español.

A diferencia de los grandes repositorios internacionales, concebidos
principalmente para clasificación supervisada, DermapixelAI~1.0 ha sido
diseñado desde su origen para soportar tareas heterogéneas
---clasificación, recuperación multimodal, razonamiento asistido por RAG
y adaptación eficiente de modelos fundacionales---, aproximándose más a
las necesidades reales de un entorno clínico digital. Esta orientación lo
convierte, antes que en un \emph{dataset}, en la plataforma experimental
sobre la que se han generado los resultados de las
secciones~\ref{sec:dermapixel-results}--\ref{sec:dermapixel-extra-sintesis}.

La interpretación debe acotarse: se trata de una primera versión de
tamaño limitado y con una cola larga marcada a nivel L3, por lo que las
cifras derivadas requieren la estimación robusta discutida en
sección~\ref{sec:disc-cv-robustez} y no admiten extrapolación directa a
la práctica asistencial.

\subsection{Eficiencia en etiquetas y régimen de saturación}
\label{sec:disc-eficiencia}

Si el preentrenamiento de dominio es tan determinante, cabe
preguntarse cuántas etiquetas son realmente necesarias para
explotarlo. La curva de eficiencia en etiquetas
(sección~\ref{sec:label-eff-results}, tabla~\ref{tab:label-eff})
muestra un comportamiento característico de los modelos fundacionales:
la mayor parte de la capacidad discriminativa ya está codificada en las
representaciones aprendidas durante el preentrenamiento, de modo que el
clasificador supervisado sólo necesita una pequeña cantidad de ejemplos
para adaptarse a una tarea concreta.

Sobre HAM10000, PanDerm Large alcanza $0{,}850$ de \emph{accuracy} y
$0{,}918$ de AUROC utilizando únicamente el $1\%$ del conjunto de
entrenamiento ($82$ imágenes), lo que equivale al $95{,}7\%$ y al
$96{,}2\%$ del rendimiento obtenido con el conjunto completo. Con sólo
el $5\%$ de los datos ($410$ imágenes) la AUROC asciende a $0{,}932$,
el $97{,}7\%$ del máximo observado y a sólo $2{,}2$\,pp de la cifra con
el $100\%$. La \emph{balanced accuracy} continúa mejorando con el
incremento de muestras y no alcanza su mejor valor ($0{,}589$) hasta el
$50\%$ del conjunto, reflejando que las clases minoritarias siguen
beneficiándose de un mayor volumen de ejemplos.

El mismo patrón aparece sobre PAD-UFES-20 ($N_{\mathrm{train}} = 1\,493$
con el conjunto completo). Con únicamente $14$ imágenes de entrenamiento
($1\%$), PanDerm Large obtiene $0{,}740$ de \emph{accuracy} y una AUROC
de $0{,}907$, equivalente al $95{,}5\%$ del rendimiento alcanzado con
todas las etiquetas disponibles. La rápida saturación observada en ambos
datasets sugiere que gran parte del conocimiento visual relevante para la
dermatología ya se encuentra internalizado en las representaciones del
codificador.

Este comportamiento tiene una consecuencia práctica importante. En los
modelos supervisados clásicos, la obtención de nuevas etiquetas suele
constituir la principal vía para mejorar el rendimiento; en cambio, en
los modelos fundacionales dermatológicos la rentabilidad marginal de
cada nueva etiqueta disminuye rápidamente una vez superado un umbral
relativamente bajo de ejemplos. A partir de ese punto, la calidad del
preentrenamiento y la adecuación del dominio pueden resultar más
determinantes que el simple aumento del tamaño del conjunto etiquetado.

Desde una perspectiva clínica, el resultado es especialmente relevante
para enfermedades raras, centros con bajo volumen asistencial o
subgrupos infrarrepresentados, donde la obtención de miles de
anotaciones expertas es poco realista. Los datos sugieren que una
estrategia basada en codificadores fundacionales especializados y
cabezas ligeras puede ofrecer rendimientos cercanos al máximo alcanzable
incluso en escenarios de escasez de datos. Esta observación refuerza una
de las ideas centrales del trabajo: en dermatología, el cuello de
botella deja de ser necesariamente la disponibilidad de etiquetas y pasa
a ser, en muchos escenarios, la calidad del preentrenamiento del que se
parte.

\subsection{Sondeo lineal frente a ajuste fino}
\label{sec:disc-lp-vs-ft}

La elección entre congelar el encoder o reentrenarlo no tiene una
respuesta única, y esta subsección discute de qué depende. La
comparación directa entre sondeo lineal y ajuste fino sobre cuatro
datasets (sección~\ref{sec:ft-results},
tabla~\ref{tab:ft-panderm}) revela un patrón asimétrico con
implicaciones para el diseño de sistemas clínicos. Sobre HAM10000 y
PAD-UFES-20, datasets de tamaño intermedio o alto
($N \in [1\,493, 8\,207]$), el ajuste fino mejora la BAcc en
$+47{,}1$\,pp y $+19{,}4$\,pp respectivamente. Esto confirma la
utilidad de la receta original de PanDerm ---\emph{weighted random
sampler}, Mixup, CutMix y \emph{layer decay}--- para aprender las
clases minoritarias bajo desbalance acusado.

El comportamiento se invierte cuando el dataset es pequeño. Sobre
DDI, con $N_{\mathrm{train}} = 427$, el sondeo lineal supera al
ajuste fino en accuracy ($+7{,}3$\,pp), BAcc ($+2{,}1$\,pp),
AUROC ($+14{,}5$\,pp) y F1 ponderada ($+5{,}6$\,pp). El
descenso de AUROC ($0{,}827 \to 0{,}682$) es especialmente
diagnóstico: el encoder ajustado pierde capacidad de ranking sobre
el \emph{test}, lo que es compatible con sobreajuste a la
distribución de entrenamiento pese a las augmentaciones empleadas.
Sobre MSKCC ($N_{\mathrm{train}} = 6\,189$), la mejora de BAcc es
marginal ($+3{,}0$\,pp) y la AUROC desciende ligeramente
($-2{,}7$\,pp). La asimetría admite una explicación común: el ajuste
fino sólo rentabiliza su mayor capacidad cuando el conjunto de
entrenamiento basta para estimar los nuevos pesos sin memorizar la
distribución; por debajo de ese umbral, congelar el encoder actúa como
regularización implícita y preserva la capacidad de \emph{ranking}
heredada del preentrenamiento.

Estos tres casos delimitan un régimen operativo cuantificable. Para
datasets con $N_{\mathrm{train}} \lesssim 500$ imágenes, el sondeo
lineal sobre PanDerm Large es la opción de referencia. Para
$N_{\mathrm{train}} \gtrsim 5\,000$, el ajuste fino con la receta
original es preferible, siempre que existan clases minoritarias
clínicamente relevantes que el sondeo lineal no logre discriminar.
Entre ambos regímenes, la decisión depende del desbalance de clases
y del coste computacional asumible.

La oposición no se reduce, sin embargo, a congelar frente a ajustar el
encoder completo: cabe una tercera vía intermedia. La evidencia sobre
DermapixelAI~1.0 a nivel L2 (sección~\ref{sec:dermapixel-ft})
indica que la adaptación de bajo rango mediante LoRA mejora la BAcc frente
al ajuste fino parcial denso entrenando del orden de $48$ veces menos
parámetros. Este patrón condiciona la estrategia recomendable bajo $N$
reducido: la adaptación eficiente es preferible al ajuste denso, no sólo
por coste computacional sino por menor exposición al sobreajuste. La
implicación de diseño es directa y se materializa en la decisión
arquitectónica de Dermapixel R0 (sección~\ref{sec:dermapixel-spanderm-v0}),
que adopta LoRA sobre los últimos bloques de PanDerm Large en lugar de un
ajuste fino denso sobre el mismo encoder.

En conjunto, la evidencia ordena las tres estrategias en una regla
escalonada por volumen de datos: con $N_{\mathrm{train}} \lesssim 500$
imágenes, sondeo lineal sobre el encoder congelado; en el régimen
intermedio, adaptación de bajo rango (LoRA), preferible al ajuste denso
por su menor exposición al sobreajuste a igualdad o mejora de
rendimiento; y sólo con $N_{\mathrm{train}} \gtrsim 5\,000$ imágenes y
clases minoritarias clínicamente relevantes, ajuste fino denso con la
receta original. La capacidad del modelo a desplegar debe escalarse, por
tanto, con los datos disponibles, no fijarse a priori.

\subsection{Segmentación promptable y adaptación de bajo coste}
\label{sec:disc-segmentacion}

La segmentación aporta el hallazgo de mayor magnitud absoluta del
trabajo, y conviene discutir tanto su alcance como su origen. Sobre
ISIC2018 (sección~\ref{sec:seg-results}, tabla~\ref{tab:seg}),
SAM2.1-Large con \emph{fine-tuning} de su decodificador alcanza
Dice $0{,}947$ en el conjunto de test. Supera en $+2{,}6$\,pp la
cifra publicada por Yan et~al.~\cite{panderm2025} ($0{,}921$) y
en $+5{,}3$\,pp al baseline CAE-seg con backbone PanDerm Large
entrenado desde cero ($0{,}894$). El ajuste mantiene congelado el
\emph{image\_encoder} y entrena sólo el decodificador de máscaras y
el \emph{promptable encoder}. Por eso los pesos actualizados se
limitan a $\sim 4{,}2$\,M sobre los $\sim 227$\,M de
SAM2.1-Large (menos del $2\%$), con un coste computacional bajo.

La transferencia fuera de dominio es prácticamente nula: el Dice apenas
varía sobre ISIC2017 y PH2 respecto a ISIC2018. Este patrón sugiere que
las representaciones del encoder visual congelado (Hiera-Large) codifican
primitivas geométricas y de contorno suficientemente generales, de modo
que el ajuste del decodificador de máscaras y del \emph{promptable
encoder} captura los matices del dominio dermatoscópico sin sobreajustar a
una fuente de adquisición concreta. La magnitud del efecto de la
adaptación de dominio, cuantificada en sección~\ref{sec:seg-results}
frente a SAM2 \emph{zero-shot}, confirma que el coste de adaptación se
concentra en el decodificador y no en el encoder.

La comparativa pareada de arquitecturas promptables
(sección~\ref{sec:seg-arquitecturas},
tabla~\ref{tab:seg-arquitecturas}) matiza, no obstante, dónde reside
ese valor. La generación de la arquitectura pesa más que el
preentrenamiento médico: SAM2.1-tiny generalista supera a MedSAM~v1
con preentrenamiento médico, y el tamaño de \emph{backbone} no
aporta en el régimen \emph{box-prompted}. La especialización
dermatológica del segmentador parece, por tanto, un factor
secundario frente al diseño del modelo. Esta asimetría contrasta con
el peso decisivo del preentrenamiento de dominio observado en
clasificación (sección~\ref{sec:disc-especializacion}) e invita a
tratar por separado ambas familias de tareas.

Este resultado sugiere que la segmentación constituye un caso
particular dentro de los modelos fundacionales médicos. Mientras que
en clasificación la especialización de dominio emerge como el principal
determinante del rendimiento (sección~\ref{sec:disc-especializacion}),
en segmentación la capacidad de la arquitectura parece dominar sobre el
origen del preentrenamiento. Dicho de otro modo, reconocer \emph{qué}
lesión aparece en una imagen sigue dependiendo en gran medida del
conocimiento dermatológico adquirido durante el preentrenamiento,
mientras que delimitar \emph{dónde} termina la lesión parece depender
principalmente de mecanismos geométricos y espaciales más generales.

El control automático sin caja
(sección~\ref{sec:seg-control-automatico}) reordena, sin embargo,
esta jerarquía de factores. Si las cinco arquitecturas
\emph{promptables} caben en $\sim 1$\,pp pero retirar el \emph{prompt}
cuesta $7{,}5$\,pp, el determinante real del rendimiento no es el
segmentador sino la caja: la información de localización pesa entre
seis y siete veces más que la elección de modelo. La consecuencia de
diseño es directa ---el esfuerzo de ingeniería rinde más en el
detector que produce la caja que en sustituir un segmentador por
otro--- y conecta con la decisión de producto de ofrecer un modo
interactivo, en el que el clínico aporta la caja, junto a uno
automático. Desde una perspectiva metodológica, este hallazgo obliga
además a reinterpretar las comparaciones habituales entre segmentadores:
cuando la diferencia entre arquitecturas es del orden de un punto
porcentual y la calidad del \emph{prompt} modifica el rendimiento en más
de siete, una parte sustancial de la literatura comparativa puede estar
midiendo indirectamente diferencias en localización más que diferencias
reales entre segmentadores. Para este último, el análisis de robustez
(sección~\ref{sec:seg-robustez}) deja una especificación poco
habitual pero accionable: dado que apretar la caja degrada el Dice
mucho más que aflojarla, el detector debe sesgarse hacia cajas
ligeramente holgadas. La validación cruzada de cinco pliegues sitúa
además el \emph{headline} en $0{,}9554\pm 0{,}0010$, lo que eleva la
cifra de segmentación de estimación puntual a estimación con
intervalo, a diferencia del resto de protocolos del trabajo.

Conviene, por último, no sobreinterpretar la magnitud absoluta. El
análisis sobre anotación múltiple
(sección~\ref{sec:seg-techo-anotador}) muestra que el acuerdo entre
dermatólogos sobre el mismo contorno es de apenas $0{,}728$ de Dice,
mientras que el modelo concuerda con el consenso en $0{,}915$. El
$0{,}9556$ se sitúa, por tanto, cerca del techo que impone la propia
ambigüedad del borde lesional: las diferencias de $\sim 1$\,pp entre
arquitecturas son inferiores al desacuerdo inter-anotador y el «error»
residual frente a una referencia única es, en buena medida, ruido de
etiqueta. La salvedad pertinente es que esa consistencia convive con
una ligera sobreconfianza en los bordes dudosos, donde el modelo traza
una frontera nítida que los humanos no comparten. En consecuencia, el
valor de Dice debe interpretarse como una aproximación al consenso
humano más que como una medida absoluta de verdad geométrica: una vez
alcanzado el entorno de $0{,}95$ frente a una referencia única, el
margen de mejora técnica disponible pasa a estar condicionado por la
propia incertidumbre de la anotación clínica.

\subsection{Sensibilidad al tokenizador en zero-shot multimodal}
\label{sec:disc-zs}

El rendimiento zero-shot no depende sólo del modelo, sino de cómo se
le interroga. La evaluación sobre HAM10000
(sección~\ref{sec:zs-results}, tabla~\ref{tab:zs-ham}) revela una
sensibilidad acusada al tokenizador y a la formulación de los
\emph{prompts}. La AUROC oscila entre $0{,}366$ (DermLIP v1/v2 con
tokenizador PubMedBERT y \emph{prompts} genéricos) y $0{,}854$
(DermLIP original con tokenizador GPT-2/CLIP y \emph{prompts} A3 de
Derm1M). Es un rango de $48{,}8$\,pp para variaciones que no tocan
los pesos visuales, sino sólo la rama lingüística y la formulación
textual de las clases.

A diferencia de los clasificadores supervisados convencionales, donde
la etiqueta final se obtiene a partir de una cabeza entrenada
explícitamente para la tarea, los sistemas \emph{vision}-lenguaje
\emph{zero-shot} realizan una comparación geométrica entre
\emph{embeddings} visuales y textuales. En consecuencia, la calidad de
la predicción depende simultáneamente de ambos espacios y, sobre todo,
de su alineamiento mutuo: un encoder visual excelente puede producir
resultados mediocres si las representaciones textuales no ocupan la
región adecuada del espacio compartido.

El valor AUROC $<0{,}5$ de la configuración inicial es diagnóstico de un
desalineamiento entre el espacio textual proyectado por PubMedBERT y el
espacio visual del encoder de DermLIP v1/v2, hasta el punto de invertir la
correlación respecto a la tarea de clasificación. El experimento de
modificaciones aisladas (sección~\ref{sec:zs-results}) permite atribuir la
recuperación del rendimiento a la combinación de dos factores ---el cambio
de tokenizador y el cambio al conjunto de \emph{prompts} A3---, sin que la
evidencia disponible permita aislar una causa única de forma concluyente,
puesto que ambos operan sobre la misma rama textual. La lectura no debe,
por tanto, formularse como causalidad absoluta de un único componente,
sino como el efecto conjunto del tokenizador, los \emph{prompts} y el
alineamiento texto-imagen resultante. Los resultados sugieren que, en
régimen \emph{zero-shot}, la rama textual deja de ser un mecanismo de
entrada y pasa a formar parte efectiva del modelo de decisión: desde
esta perspectiva, el tokenizador, los \emph{prompts} y el encoder
textual no son hiperparámetros externos, sino componentes funcionales
del clasificador. De aquí se sigue una implicación práctica directa: todo despliegue \emph{zero-shot} debe validar
empíricamente, sobre un conjunto del dominio, no sólo el modelo sino
también su tokenizador y la formulación de los \emph{prompts}, porque la
transferencia desde la configuración publicada por los autores no está
garantizada. El experimento cuestiona, además, una práctica habitual en
la literatura multimodal: reportar resultados \emph{zero-shot} utilizando
una única configuración de \emph{prompts}. Sobre los datos de este
trabajo, la variabilidad inducida por la formulación textual supera
ampliamente la diferencia entre algunas arquitecturas visuales, lo que
sugiere que las comparaciones entre modelos sin controlar el espacio de
\emph{prompts} pueden resultar difíciles de interpretar.

La evaluación adicional sobre DermapixelAI~1.0
(sección~\ref{sec:dermapixel-zs}, tabla~\ref{tab:dermapixel-zs})
extiende esta lectura al idioma. DermLIP v2 con \emph{prompts} en
inglés alcanza Acc@1 $= 0{,}500$ sobre la categoría etiológica
L1, frente a $0{,}361$ con \emph{prompts} equivalentes en
castellano ($-13{,}9$\,pp); SigLIP-SO400M cae $-11{,}1$\,pp en
la misma comparación. BiomedCLIP, en cambio, mantiene Acc@1
$= 0{,}389$ en ambos idiomas. Esto sitúa la penalización por
castellano como dependiente del dataset textual de preentrenamiento.
La magnitud de la caída motiva, como línea futura inmediata,
construir modelos vision-lenguaje contrastivos con rama textual
entrenada sobre material clínico en castellano. Esta penalización
proporciona, además, una posible explicación mecanística para parte de
los resultados obtenidos posteriormente con Dermapixel R0
(sección~\ref{sec:disc-especializacion}): si una simple traducción del
vocabulario diagnóstico reduce el rendimiento entre $11$ y $14$\,pp,
resulta razonable esperar que una adaptación explícita de la rama
textual al lenguaje clínico castellano aporte beneficios incluso cuando
la representación visual permanece inalterada.

\subsection{Brecha entre modelos específicos y LLMs multimodales}
\label{sec:disc-llm}

La sensibilidad de la rama contrastiva al texto invitaba a
comprobar si los grandes modelos de lenguaje multimodales, con
dataset de preentrenamiento mucho mayores, cierran la distancia con
los especialistas. La comparación de paradigmas sobre HAM10000
(sección~\ref{sec:llm-results}, tabla~\ref{tab:llm-comparativa})
muestra lo contrario y fija un orden de magnitud entre familias
relevante para elegir arquitectura en escenarios clínicos. El
gradiente de accuracy entre el mejor especialista ajustado (PanDerm
Base con ajuste fino y \emph{test-time augmentation}) y el mejor LLM
multimodal \emph{zero-shot} (GPT-4o) supera los $43$\,pp; la
tabla~\ref{tab:llm-comparativa} desglosa la escala completa de
configuraciones intermedias. La magnitud de la diferencia observada
sugiere que, al menos en dermatología visual, el conocimiento específico
del dominio continúa siendo más determinante que la escala general del
modelo: la disponibilidad de miles de millones de parámetros y de
grandes capacidades lingüísticas no compensa la ausencia de una
representación dermatológica especializada.

Tres observaciones cualitativas, detalladas en
sección~\ref{sec:llm-results}, matizan esta brecha sin reducirla.
Primero, la adaptación con LoRA sobre MedGemma~27B mejora la accuracy
sobre los datasets vistos en entrenamiento pero no transfiere de forma
sistemática al resto, lo que aconseja reportar su rendimiento por dataset
y no agregado. Segundo, GPT-4o y GPT-4o-mini exhiben un sesgo de error
hacia clases minoritarias del \emph{prompt} ---con sobreproducción de
predicciones de dermatofibroma y melanoma muy por encima de su prevalencia
real---, lo que sugiere que responden al espacio léxico de la consigna
tanto como a la evidencia visual. Este patrón es compatible con un
razonamiento guiado parcialmente por asociaciones semánticas aprendidas
durante el preentrenamiento generalista y no exclusivamente por evidencia
visual, por lo que la interpretación de las respuestas de los LLM
multimodales requiere cautela cuando se aplican a taxonomías clínicas
especializadas. Tercero, BLIP-2 e InstructBLIP con
\emph{Q-Former} preentrenado producen una accuracy incompatible con
cualquier uso clínico, atribuible a la falta de adaptación al vocabulario
dermatológico. Las tres observaciones convergen en que la distancia entre
especialistas y LLMs multimodales generalistas no se explica por capacidad
bruta del modelo, sino por su grado de adaptación al dominio. La
conclusión queda acotada, en todo caso, a los \emph{prompts}
estandarizados y a la versión de los modelos cerrados disponible en la
fecha de evaluación (sección~\ref{sec:lim-pesos}): no afirma el techo
teórico de estos modelos, sino su rendimiento en el régimen de uso
evaluado.

De forma consistente con otros experimentos del trabajo, la escala del
modelo no emerge como la variable dominante. En segmentación el principal
determinante fue la disponibilidad del \emph{prompt} espacial
(sección~\ref{sec:disc-segmentacion}); en \emph{zero-shot} multimodal, la
alineación texto-imagen (sección~\ref{sec:disc-zs}); y en clasificación
dermatológica, el preentrenamiento de dominio
(sección~\ref{sec:disc-especializacion}). La capacidad paramétrica por sí
sola no explica el rendimiento observado en ninguno de los tres regímenes.

La brecha tiene además una vertiente de coste. La evaluación con
APIs comerciales sobre los tres datasets cuesta unos \$4{,}07.
MedGemma~27B, en cambio, opera localmente en BF16 sobre los
$128$\,GB de memoria unificada de la DGX Spark sin cuantización,
con latencia media de $1{,}49$ segundos por imagen. La diferencia
operativa importa para un despliegue hospitalario: los modelos
cerrados añaden coste recurrente y dependencia de un proveedor
externo, sin mejora apreciable de rendimiento clínico. Desde una
perspectiva hospitalaria, el resultado desplaza el criterio de selección
desde la popularidad del modelo hacia su adecuación al dominio: bajo las
condiciones evaluadas, los modelos fundacionales dermatológicos ofrecen
simultáneamente mejor rendimiento, menor coste operativo y mayor
soberanía tecnológica que las alternativas multimodales cerradas de
propósito general.

\subsection{Equidad por fototipo cutáneo}
\label{sec:disc-equidad}

La generalización entre fototipos es una condición de seguridad
clínica, y los resultados la cuantifican sin ambigüedad. La
evaluación sobre Fitzpatrick17k completo
(sección~\ref{sec:equidad-results}, tablas~\ref{tab:equidad-fototipo}
y~\ref{tab:cross-eval}) reproduce un patrón conocido en la
literatura: los cinco encoders contrastados presentan \emph{gap}~I-VI
positivo en la formulación de tres clases de malignidad, con valores
entre $+0{,}124$ y $+0{,}306$\,pp de BAcc. Ese \emph{gap} no se
interpreta en aislamiento, sino junto a la BAcc global de cada
modelo. DermLIP v2 obtiene la mejor BAcc global ($0{,}721$) y la
mejor AUROC global ($0{,}909$), pero su \emph{gap}~I-VI es
$+0{,}281$\,pp: el rendimiento sobre FST~VI ($0{,}434$) cae
mucho respecto a FST~III ($0{,}765$). PanDerm Large presenta un
\emph{gap} menor ($+0{,}217$\,pp) con BAcc global ligeramente
inferior ($0{,}699$). BiomedCLIP exhibe el \emph{gap} mínimo
($+0{,}124$\,pp), pero la peor BAcc global ($0{,}590$).

Esta última asimetría ---\emph{gap} mínimo con BAcc global
inferior--- ilustra que la métrica de equidad aislada no es
decisional. \textbf{Un modelo uniformemente malo entre subgrupos puede
mostrar baja desigualdad aritmética sin ser equitativo en sentido
sustantivo}, porque su rendimiento es bajo en todos los fototipos por
igual. El \emph{gap} mínimo de BiomedCLIP no señala, por tanto, un
modelo más justo, sino uno uniformemente peor: la equidad no puede
medirse al margen del rendimiento que reparte. Por eso BAcc global y
\emph{gap} deben leerse de forma conjunta. Considerando conjuntamente
rendimiento global y magnitud del \emph{gap}, PanDerm Large presenta el
equilibrio más favorable entre ambas dimensiones dentro de los modelos
evaluados, seguido de DermLIP v2 cuando se ponderan el rendimiento sobre
FST~I-IV y la AUROC global.

Desde una perspectiva asistencial, el hallazgo implica que un modelo
validado únicamente mediante métricas globales puede ocultar pérdidas
relevantes de sensibilidad en pacientes con fototipos menos
representados. La validación estratificada por fototipo debería
considerarse, por tanto, un requisito de seguridad y no únicamente una
evaluación complementaria de equidad. Esta exigencia adquiere además
relevancia regulatoria: los marcos emergentes de evaluación de sistemas
de IA sanitaria enfatizan la necesidad de identificar degradaciones
sistemáticas del rendimiento entre subpoblaciones, especialmente cuando
dichas diferencias pueden traducirse en desigualdades diagnósticas.

El experimento cruzado \emph{Light}~$\leftrightarrow$~\emph{Dark}
(tabla~\ref{tab:cross-eval}) añade una segunda lectura. La caída de
rendimiento al entrenar sólo con fototipos claros y evaluar sobre
fototipos oscuros (L$\to$D vs L$\to$L) es máxima en DINOv2
($\Delta = -0{,}132$\,pp BAcc) y mínima en BiomedCLIP
($\Delta = -0{,}063$\,pp). DermLIP v2 mantiene la menor caída
Dark$\to$Light vs Dark$\to$Dark ($\Delta = +0{,}016$\,pp), lo que
sugiere mayor invariancia de sus representaciones al fototipo. El
patrón es coherente con su dataset de preentrenamiento Derm1M,
derivado de fuentes web abiertas con representación variable de
fototipos.

Este experimento cruzado es, en cierto sentido, más informativo que el
\emph{gap}~I-VI descriptivo, porque responde a una pregunta causal:
qué ocurre al entrenar con unos fototipos y desplegar sobre otros. La
caída sistemática observada en el escenario L$\to$D confirma que el
sesgo no se limita a diferencias de prevalencia entre grupos, sino que
emerge incluso cuando el entrenamiento excluye explícitamente parte de
la variabilidad fenotípica. El fenómeno es compatible con una
dependencia real de las representaciones visuales respecto al tono
cutáneo, y no con un mero artefacto de muestreo.

Ninguno de los cinco encoders evaluados elimina completamente la
degradación observada en fototipos elevados. La convergencia de
resultados entre arquitecturas muy distintas sugiere que la limitación
no reside exclusivamente en un modelo concreto, sino que refleja un
problema más amplio de representación de fototipos oscuros en los datos
disponibles para entrenamiento. Los resultados justifican, en
consecuencia, la incorporación explícita de diversidad fenotípica en
futuras versiones de DermapixelAI: la expansión del repositorio clínico
del HUSLL (sección~\ref{sec:disc-dermapixel-contrib}) ofrece una
oportunidad para aumentar la representación de fototipos
infrarrepresentados y reducir la dependencia de distribuciones
procedentes de datasets públicos.

\subsection{El cuello de botella ya no es sólo el modelo}
\label{sec:disc-cuello-botella}

La evidencia acumulada a lo largo del trabajo sugiere que el rendimiento
dermatológico no está limitado principalmente por la arquitectura del
modelo. En varios experimentos, modificaciones en el diseño del sistema
---tokenización, \emph{prompts}, estrategia de adaptación, estructura
ontológica, equilibrio de clases, recuperación por similitud o
validación por subgrupos--- producen variaciones comparables o
superiores a las obtenidas al sustituir el codificador visual. Estas
decisiones, más que la elección aislada de un codificador, determinan el
rendimiento alcanzable, y la lectura desplaza parcialmente el foco desde
el diseño del modelo hacia el diseño del sistema completo.

Esta evolución es coherente con la madurez progresiva de los modelos
fundacionales. A medida que las representaciones visuales alcanzan
niveles elevados de calidad, el margen de mejora se desplaza desde la
arquitectura hacia aspectos tradicionalmente considerados periféricos:
datos, ontologías, alineamiento multimodal, calibración, equidad e
integración operativa. En este sentido, Dermapixel R0 no representa la
construcción de un modelo mayor, sino de un sistema mejor adaptado al
problema: la contribución no reside en la escala, sino en identificar
qué componentes limitan realmente el rendimiento y actuar sobre ellos.

Desde una perspectiva de despliegue clínico, la consecuencia es clara:
optimizar exclusivamente el modelo puede producir retornos decrecientes
si no se acompaña de mejoras equivalentes en los datos, la taxonomía
diagnóstica, la estrategia de adaptación, la validación por subgrupos y
la integración en el flujo asistencial. Esta lectura constituye el
objeto de las implicaciones operativas que se desarrollan en la
sección~\ref{sec:implicaciones}.

\subsection{Estimación robusta del rendimiento sobre datasets de tamaño moderado}
\label{sec:disc-cv-robustez}

La fiabilidad de las cifras sobre un dataset pequeño depende del método de
estimación empleado. La validación cruzada de cinco \emph{folds} agrupada
por caso clínico sobre DermapixelAI~1.0
(sección~\ref{sec:dermapixel-results}, tabla~\ref{tab:dermapixel-cv})
ilustra un punto metodológico relevante para evaluar modelos
fundacionales sobre datasets de tamaño moderado: la diferencia entre la
cifra puntual del \emph{split} fijo y la media de los \emph{folds} es
sustancial sobre las métricas dependientes de umbral (BAcc) y casi nula
sobre la AUROC \emph{one-vs-rest}. Esta observación matiza la lectura de
toda cifra de BAcc puntual sobre un \emph{test} reducido. La validación
cruzada transforma así la interpretación del resultado desde una
estimación puntual hacia una estimación de su variabilidad: la cifra
puntual no es el rendimiento, sino una observación; el rendimiento es una
distribución. En datasets clínicos de tamaño moderado, esta información
resulta tan importante como la propia métrica central. De aquí se siguen
dos lecturas operativas. \textit{Primero}, la cifra del \emph{split}
publicado conserva validez como punto operativo de referencia ---y como
base de la comparación cruzada con la evaluación \emph{zero-shot} sobre el
mismo \emph{test} (sección~\ref{sec:dermapixel-zs})---, pero requiere la
estimación cruzada para defender el rendimiento esperado del sistema.
\textit{Segundo}, sobre datasets clínicos de tamaño moderado la AUROC es la
métrica primaria recomendable para comparar configuraciones de
\emph{transfer learning} congelado, mientras que la BAcc y la
\emph{accuracy} deben acompañarse de su intervalo de confianza o de su
estimación cruzada por \emph{folds}. La estabilidad de la AUROC frente a
las métricas dependientes de umbral explica además su uso recurrente a lo
largo del trabajo como métrica primaria de comparación entre
codificadores: mientras que la BAcc depende de una frontera de decisión
concreta, la AUROC evalúa la capacidad de ordenación del modelo y resulta
menos sensible a fluctuaciones muestrales en conjuntos reducidos.

En consecuencia, los \emph{rankings} observados entre modelos sobre
DermapixelAI~1.0 deben interpretarse como estimaciones sujetas a
incertidumbre muestral y no como jerarquías absolutas. La validación
cruzada agrupada por caso proporciona una aproximación más robusta al
rendimiento esperado en despliegue que cualquier partición fija
individual, y protege la lectura de las cifras de los niveles L2 y L3
---los de \emph{test} más reducido--- frente a la sobreinterpretación de
diferencias dentro del margen de error.

Esta cautela enmarca también la lectura de los tres \emph{endpoints} de
malignidad contrastados formalmente con el test de DeLong
(sección~\ref{sec:equidad-delong}). El cierre estadístico se concentra
en Fitzpatrick17k ---el único con potencia muestral suficiente, $226$
positivos sobre $1\,658$---, donde los dos codificadores específicos de
dermatología (PanDerm Large y DermLIP v2) empatan entre sí
($\Delta\mathrm{AUROC} = 0{,}0036$, $p = 0{,}55$) y superan de forma
significativa a los generalistas y a PanDerm Base. Sobre los dos
\emph{endpoints} de menor tamaño ---HAM10000 ($70$ melanomas) y
DermapixelAI~1.0 ($58$ malignos en validación cruzada)--- ningún
contraste alcanza significación tras la corrección de Holm: las
diferencias de AUROC son pequeñas ($\leq 0{,}006$ en HAM,
$\leq 0{,}05$ en DermapixelAI) y los intervalos se solapan ampliamente,
un patrón compatible con la falta de potencia y no con la ausencia de
efecto. El titular sostenible es, por tanto, la jerarquía «específico de
dominio $>$ generalista», con cierre estadístico en Fitzpatrick17k y
corroboración descriptiva en los otros dos; el contraste no respalda una
ordenación fina entre los dos líderes específicos. Reportar los tres
\emph{endpoints} con su corrección por comparaciones múltiples ---y
declarar cuál tiene potencia y cuál no--- es la práctica que las guías
de reporte para inteligencia artificial en imagen médica (TRIPOD+AI)
prescriben, y que distingue una afirmación de superioridad
estadísticamente sostenida de una mera diferencia puntual.

\section{Comparación con la literatura}
\label{sec:lit}

Los resultados discutidos en la sección anterior adquieren significado
únicamente en relación con la evidencia previa disponible. Por ello, esta
sección contrasta los hallazgos obtenidos con los principales trabajos de
referencia en modelos fundacionales dermatológicos, segmentación médica,
aprendizaje multimodal, equidad por fototipo e interpretabilidad basada
en conceptos.

El objetivo no es únicamente verificar la reproducibilidad de resultados
previamente publicados, sino también identificar los puntos de
convergencia y divergencia que emergen al evaluar estos modelos sobre
datasets externos, en castellano y bajo condiciones de despliegue
diferentes a las descritas en los trabajos originales. Este contraste
permite situar las contribuciones de DermapixelAI~1.0 y Dermapixel R0
dentro del estado actual de la investigación, delimitando qué
observaciones confirman la literatura existente, cuáles la matizan y
cuáles aportan evidencia nueva.

\subsection{Reproducción del benchmark de PanDerm}
\label{sec:lit-panderm}

El primer contraste con la literatura corresponde a la
reproducibilidad del propio modelo fundacional.
Yan et~al.~\cite{panderm2025} reportan para PanDerm Large una mejora
media del $+10{,}2\%$ sobre el estado del arte previo a través de los
$28$ \emph{downstream tasks} de su benchmark interno. Las cifras de
sondeo lineal obtenidas en este trabajo sobre los datasets compartidos
(HAM10000, PAD-UFES-20 y Derm7pt; tabla~\ref{tab:lp-panderm}) se sitúan
dentro del rango descrito por los autores y reproducen cualitativamente
las relaciones de rendimiento entre PanDerm Base y PanDerm Large.

Esta concordancia resulta especialmente relevante porque la evaluación
se ha realizado bajo condiciones de ejecución distintas a las del
trabajo original. Los experimentos se ejecutaron sobre arquitectura
ARM64 (\texttt{aarch64}), aceleración CUDA~13.0 y precisión BF16 en la
plataforma NVIDIA Grace Hopper GB10, mientras que la validación original
se realizó sobre infraestructura x86\_64. La proximidad de los
resultados ---dentro de las tolerancias esperables por aritmética BF16 y
por diferencias de versiones de \texttt{cuDNN}--- indica que el
comportamiento del modelo es robusto frente a diferencias razonables de
hardware y entorno de ejecución, y documenta la portabilidad práctica
del \emph{pipeline} hacia arquitecturas de nueva generación orientadas a
IA.

La magnitud de la mejora observada también es consistente con la
literatura. Sobre los diez datasets externos evaluados en este trabajo,
PanDerm Large supera a PanDerm Base en $+4{,}1$\,pp de \emph{accuracy} y
$+9{,}0$\,pp de \emph{balanced accuracy}. Este último valor se aproxima
notablemente al $+10{,}2\%$ reportado por Yan et~al. sobre su benchmark
interno, pese a que los datasets empleados aquí son externos al
desarrollo original. La convergencia entre ambos resultados sugiere que
la ventaja de PanDerm Large no depende exclusivamente de los conjuntos
utilizados durante el desarrollo del modelo, sino que se mantiene cuando
se evalúa sobre cohortes independientes.

Más allá de validar la reproducibilidad numérica, este resultado aporta
evidencia externa sobre la generalización del modelo fundacional. En un
contexto donde numerosos modelos recientes muestran pérdidas apreciables
de rendimiento al abandonar los datasets de desarrollo, la estabilidad
observada entre el benchmark interno del grupo de Monash y los conjuntos
públicos utilizados en este trabajo constituye un indicador favorable de
robustez y transferibilidad.

\subsection{Segmentación binaria de lesión}
\label{sec:lit-seg}

En segmentación, la cifra obtenida también supera la referencia
publicada y conviene interpretar las razones probables de esta
diferencia. SAM2.1-Large con \emph{fine-tuning} del decodificador sobre
ISIC2018 alcanza Dice $= 0{,}947$, mejorando en $+2{,}6$\,pp el
resultado reportado por Yan et~al.~\cite{panderm2025} ($0{,}921$ Dice
con $100$ épocas).

Tres factores pueden explicar plausiblemente esta diferencia.
\textit{Primero}, la arquitectura \emph{promptable} de SAM2.1 incorpora
una señal espacial explícita mediante la \emph{bounding box}, lo que
restringe el espacio de búsqueda de la segmentación y reduce la
ambigüedad de los contornos. \textit{Segundo}, el encoder Hiera-Large de
SAM2.1, preentrenado sobre $\sim 11$ millones de imágenes y más de mil
millones de máscaras, proporciona representaciones geométricas y de
contorno extremadamente robustas, probablemente más adecuadas para esta
tarea que las aprendidas por encoders entrenados exclusivamente sobre
imágenes dermatológicas. \textit{Tercero}, la adaptación restringida al
decodificador de máscaras y al \emph{promptable encoder} concentra la
capacidad de ajuste en los componentes directamente responsables de la
segmentación, minimizando el riesgo de sobreajuste y reduciendo
significativamente el coste computacional frente a un \emph{fine-tuning}
más amplio.

Los experimentos adicionales realizados en este trabajo refuerzan esta
interpretación. La comparación pareada entre arquitecturas
\emph{promptables} (sección~\ref{sec:seg-arquitecturas}) muestra que las
diferencias entre segmentadores son relativamente pequeñas ($\approx 1$
\,pp de Dice), mientras que retirar la caja-\emph{prompt} produce una
degradación de $7{,}5$\,pp (sección~\ref{sec:seg-control-automatico}).
El principal determinante del rendimiento no parece ser, por tanto, el
segmentador concreto, sino la disponibilidad de una localización inicial
fiable de la lesión.

La literatura reciente ha tendido a interpretar la mejora de la
segmentación médica principalmente como consecuencia de arquitecturas
cada vez mayores o de preentrenamientos más especializados. Los
resultados obtenidos aquí sugieren una lectura parcialmente distinta:
una vez alcanzado un determinado nivel de calidad en las
representaciones visuales, el beneficio marginal de aumentar el tamaño
del modelo es reducido frente al efecto de disponer de información
espacial explícita y de una adaptación específica de bajo coste.

La validación cruzada (Dice $= 0{,}9554 \pm 0{,}0010$) y el análisis
inter-anotador aportan además una lectura complementaria. El acuerdo
medio entre dermatólogos sobre el contorno de una misma lesión
($0{,}728$ Dice) es sustancialmente inferior al acuerdo del modelo
respecto al consenso ($0{,}915$), lo que indica que gran parte del error
residual se sitúa ya en la incertidumbre inherente a la anotación
humana. En este contexto, la segmentación binaria de lesión
dermatoscópica parece aproximarse a un régimen de saturación práctica
cuando se combina un encoder \emph{promptable} preentrenado a gran
escala con una adaptación específica de bajo coste sobre ISIC2018.

\subsection{Comportamiento zero-shot}
\label{sec:lit-zs}

El comportamiento zero-shot observado en este trabajo es consistente
con la literatura previa sobre modelos contrastivos dermatológicos,
pero introduce una matización relevante sobre los factores que
condicionan su rendimiento. DermLIP alcanza valores de AUROC
comparables a los reportados en Derm1M~\cite{derm1m} ---máxima de
$0{,}854$ sobre HAM10000--- cuando se emplea una configuración textual
alineada con la utilizada durante su desarrollo, lo que respalda la
reproducibilidad general del paradigma.

Sin embargo, los experimentos realizados muestran que una parte
sustancial del rendimiento zero-shot no depende exclusivamente del
encoder visual ni del volumen de imágenes empleado durante el
preentrenamiento. La sensibilidad observada frente a cambios en el
tokenizador y en la formulación de los \emph{prompts}
(sección~\ref{sec:disc-zs}) ---hasta $48{,}8$\,pp de AUROC sobre el
mismo encoder visual--- sugiere que la alineación texto-imagen
constituye un componente crítico del sistema. Esta observación resulta
especialmente relevante porque la mayoría de trabajos reportan una única
configuración textual, dificultando estimar hasta qué punto el
rendimiento atribuido al modelo depende realmente de la arquitectura o
de la formulación lingüística empleada; su cuantificación sistemática no
figura en la literatura previa sobre la familia DermLIP y constituye una
aportación empírica de este trabajo.

Los resultados obtenidos son coherentes con la hipótesis de que, una vez
alcanzado un determinado nivel de calidad en las representaciones
visuales, las limitaciones del paradigma zero-shot pueden desplazarse
hacia la rama textual. La penalización observada posteriormente al
utilizar castellano en lugar de inglés sobre DermapixelAI~1.0
(sección~\ref{sec:disc-zs}) refuerza esta interpretación y sugiere que
la disponibilidad de recursos lingüísticos específicos del dominio puede
ser tan importante como la disponibilidad de imágenes.

Desde una perspectiva metodológica, estos hallazgos recomiendan cautela
al comparar modelos multimodales exclusivamente a partir de métricas
zero-shot. Diferencias atribuidas aparentemente a la arquitectura pueden
reflejar, al menos en parte, diferencias en la formulación textual
utilizada durante la evaluación. La comparación entre modelos debería
considerar explícitamente este factor, especialmente en dominios
clínicos donde la terminología es altamente especializada.

\subsection{Disparidad por fototipo}
\label{sec:lit-equidad}

La disparidad observada por fototipo es coherente con los resultados
previamente descritos en Fitzpatrick17k~\cite{fitzpatrick17k} y con la
motivación que impulsó la creación del dataset DDI~\cite{ddi}. Al igual
que en esos trabajos, los modelos evaluados muestran un descenso
sistemático del rendimiento sobre los fototipos más oscuros, lo que
confirma que el sesgo asociado a la representación desigual de tonos de
piel continúa siendo un problema abierto incluso en los modelos
fundacionales dermatológicos más recientes.

La aportación de este trabajo no reside tanto en identificar nuevamente
la existencia de esa disparidad como en compararla de forma homogénea
entre cinco familias de codificadores bajo un mismo protocolo
experimental. Esta comparación permite observar que la relación entre
rendimiento global y equidad no es necesariamente directa. DermLIP v2
obtiene la mejor capacidad discriminativa global, mientras que PanDerm
Large presenta un compromiso más equilibrado entre rendimiento y
estabilidad entre fototipos. Por el contrario, BiomedCLIP exhibe el
menor \emph{gap} entre los grupos extremos, pero también el peor
rendimiento global.

Esta última observación resulta especialmente relevante desde una
perspectiva metodológica. Un \emph{gap} reducido puede reflejar una
distribución más homogénea del error entre subgrupos, pero también puede
aparecer cuando el modelo presenta un rendimiento bajo de forma
generalizada. La denominada paradoja de BiomedCLIP pone de manifiesto
que la interpretación aislada de métricas de equidad puede conducir a
conclusiones engañosas: un modelo uniformemente mediocre puede parecer
más equitativo que otro claramente superior simplemente porque falla de
manera similar en todos los grupos.

Los resultados respaldan, por tanto, una recomendación alineada con las
discusiones recientes sobre evaluación responsable de sistemas de IA:
las métricas de equidad deben interpretarse conjuntamente con las
métricas de rendimiento global. Ni el \emph{gap} entre subgrupos ni la
\emph{accuracy} o la AUROC son suficientes por sí solas para caracterizar
el comportamiento de un sistema clínico. La evaluación de modelos
dermatológicos debería reportar simultáneamente rendimiento global,
rendimiento por subgrupo y magnitud de las diferencias observadas entre
ellos.

La evaluación cruzada \emph{Light}~$\leftrightarrow$~\emph{Dark}
(tabla~\ref{tab:cross-eval}) aporta además una segunda evidencia
consistente con la literatura reciente. Los modelos preentrenados
específicamente sobre dermatología muestran una mayor robustez frente a
cambios en la distribución de fototipos que los encoders visuales
generalistas. Este resultado sugiere que parte del beneficio del
preentrenamiento de dominio no se limita a mejorar la discriminación
diagnóstica, sino que también puede contribuir a una representación más
estable frente a variaciones demográficas clínicamente relevantes. En
conjunto, la mejora del rendimiento global no garantiza automáticamente
una mejora equivalente en equidad, y la reducción del \emph{gap} entre
subgrupos tampoco garantiza un comportamiento clínicamente adecuado
---una tensión que se retoma entre las limitaciones del trabajo
(capítulo~\ref{cap:limitaciones})---.

\section{Implicaciones operativas}
\label{sec:implicaciones}

Los resultados obtenidos no sólo permiten comparar modelos, sino también
extraer consecuencias prácticas para el diseño de futuros sistemas
dermatológicos asistidos por inteligencia artificial. A diferencia de la
literatura inicial sobre modelos fundacionales, donde el objetivo
principal era identificar un único modelo capaz de superar al estado del
arte en múltiples tareas, la evidencia acumulada en este trabajo apunta
hacia una conclusión diferente: distintos problemas dermatológicos
favorecen arquitecturas distintas.

La clasificación jerárquica, la segmentación, la recuperación de casos
similares, la búsqueda multimodal mediante lenguaje natural, la
interpretación conceptual o el cribado de seguridad presentan requisitos
técnicos diferentes, y no existe un único modelo que optimice
simultáneamente todos ellos. Los resultados sugieren que la estrategia
más defendible para un entorno clínico consiste en combinar varios
componentes especializados dentro de una arquitectura modular,
seleccionando cada uno de ellos según la tarea concreta que deba
resolver. Las subsecciones siguientes traducen los hallazgos
experimentales del capítulo~\ref{cap:resultados} en decisiones de diseño
aplicables a sistemas reales de apoyo al diagnóstico dermatológico.

\subsection{Arquitectura defendible para sistemas de apoyo al diagnóstico}
\label{sec:disc-arq}

Los cuatro bloques de evidencia anteriores convergen en una
arquitectura técnicamente defendible para sistemas de apoyo al
diagnóstico dermatológico construidos sobre los modelos
fundacionales disponibles. La estructura recomendable es la
siguiente. El encoder es PanDerm Large preentrenado, congelado por
defecto y ajustado sólo bajo demanda. La cabeza es una regresión
logística multi-clase (L-BFGS, $C=1{,}0$,
\texttt{max\_iter}~$=5\,000$) cuando el dataset de entrenamiento es
reducido ($N_{\mathrm{train}} \lesssim 500$), o una cabeza
ajustada con la receta de PanDerm Large (Mixup, CutMix,
\emph{weighted random sampler}, $50$ épocas) cuando
$N_{\mathrm{train}} \gtrsim 5\,000$ y existe desbalance acusado.
La segmentación se resuelve con SAM2.1-Large adaptado mediante
\emph{fine-tuning} del decodificador cuando la tarea requiera
localizar la lesión.

Las ramas textuales se incorporan con cautela. La rama lingüística
contrastiva (DermLIP) sólo entra con prompts validados y tokenizador
compatible, verificado sobre un conjunto de validación del dominio.
La rama LLM multimodal generalista (GPT-4o, Gemini, MedGemma) no se
usa como clasificador principal, dado el orden de magnitud de
rendimiento inferior cuantificado en sección~\ref{sec:disc-llm}.
Su rol natural es complementario: generar razonamiento textual o
explicaciones para casos ya clasificados por el encoder
especializado.

La consecuencia de diseño que se desprende del conjunto de hallazgos
(sección~\ref{sec:hallazgos}) es que la arquitectura defendible es
modular, no un modelo único: PanDerm Large para la transferencia visual
general, adaptación LoRA para la subcategoría clínica bajo $N$ reducido,
\emph{retrieval} FAISS para la cola larga, SAM2.1 para la segmentación y
un \emph{ensemble} acotado para el cribado de melanoma, seleccionando el
componente según la tarea y el nivel ontológico. Sobre esta estructura,
dos validaciones son condición necesaria antes del despliegue: la de los
\emph{prompts} y el tokenizador de cualquier rama contrastiva sobre un
conjunto del dominio, dada la sensibilidad documentada en
sección~\ref{sec:disc-zs}, y la validación por fototipo cutáneo que se
detalla en sección~\ref{sec:disc-ventana}.

Esta arquitectura modular hereda, por último, una decisión
taxonómica. La ontología no sólo permite comparar datasets, sino que
obliga a elegir a qué nivel se compara. La heterogeneidad taxonómica entre
datasets, mitigada mediante la ontología jerárquica L1/L2/L3 descrita
en sección~\ref{sec:ontologia}, sigue requiriendo una decisión
explícita sobre el nivel de agregación de los contrastes. La
comparación inter-dataset es defendible a nivel L1 (cuatro familias
etiológicas) y L2 ($43$ subcategorías clínicas) sobre el dataset
armonizado de $72\,654$ imágenes. El nivel L3 ($367$
diagnósticos individuales) conserva la cola larga propia de cada
dataset y no admite contrastes directos sin agregación adicional. La
implicación para un sistema de producción es que la jerarquía debe
mostrarse de forma explícita al usuario clínico: presentar la
predicción L3 junto a su agregación L2 y L1 da robustez frente a la
incertidumbre de la cola larga sin sacrificar especificidad cuando
ésta es fiable.

\subsection{Ventana operativa por fototipo cutáneo}
\label{sec:disc-ventana}

El rendimiento desigual entre fototipos también acota dónde puede
desplegarse el sistema con garantías. Esa ventana operativa, definida
en el experimento de equidad (tabla~\ref{tab:equidad-fototipo}),
condiciona la generalización a poblaciones con fototipos elevados.
Los cinco encoders presentan \emph{gap}~I-VI positivo entre
$+0{,}124$ y $+0{,}306$\,pp BAcc. La paradoja BiomedCLIP
---\emph{gap} mínimo asociado a la peor BAcc global--- vuelve a
ilustrar que el \emph{gap} aislado no es una métrica decisional. El
mejor compromiso entre BAcc global y \emph{gap} sobre la formulación
de tres clases es PanDerm Large, seguido de DermLIP v2 cuando se
prioriza la AUROC global. La implicación clínica es clara: desplegar
sobre poblaciones con representación significativa de fototipos
elevados exige validación específica por subgrupos y mecanismos de
revisión humana cuando el modelo opere en su régimen de menor
confianza.

\subsection{Ensemble de seguridad para detección de melanoma}
\label{sec:disc-ensemble}

La detección de melanoma admite un esquema orientado a la seguridad,
y conviene discutir tanto su valor como sus costes. El protocolo
auxiliar de ensemble (sección~\ref{sec:ensemble-results},
tabla~\ref{tab:ensemble}) muestra que la combinación OR
\emph{top-3} de PanDerm Large \emph{fine-tuned} sobre HAM10000 (M1)
y SigLIP-Large SO400M con sondeo lineal alcanza \emph{recall} de
melanoma del $100\%$ ($70$ de $70$) sobre el \emph{split} de
test de HAM10000. El análisis de errores explica por qué: los dos
melanomas que M1 no captura en \emph{top-3} sí los captura SigLIP,
lo que confirma la complementariedad de los \emph{backbones} y
justifica la elección del par. M7 (clasificador unificado sobre el
dataset armonizado de $43$ clases) no añade melanomas adicionales
sobre HAM, pero mantiene valor en escenarios multidataset por su
mayor cobertura ontológica.

Tres caveats acompañan la lectura operativa de este resultado.
(i)~El \emph{recall} del $100\%$ es \emph{top-3}, no
\emph{top-1}, por lo que exige una interfaz que presente al clínico
las tres clases más probables. (ii)~El tamaño muestral de melanomas
es $N = 70$, lo que produce intervalos de confianza al $95\%$
amplios y obliga a validación prospectiva sobre una cohorte mayor.
(iii)~El coste en falsos positivos del esquema OR \emph{top-3} es
elevado ($885$ FP sobre $1\,162$ no-melanomas): asumible en un
cribado oportunista, pero suficiente para descartar su uso aislado
como clasificación primaria. La lectura operativa, coherente con la
arquitectura modular de sección~\ref{sec:disc-arq}, es que este
\emph{ensemble} debe integrarse como mecanismo de apoyo orientado a la
sensibilidad ---priorizar la detección de melanoma asumiendo falsos
positivos--- y nunca como decisión diagnóstica autónoma.

Más allá del caso concreto del melanoma, el conjunto de resultados
obtenidos en esta tesis apunta hacia una conclusión arquitectónica más
general. Ninguno de los modelos evaluados domina simultáneamente la
clasificación jerárquica, la segmentación, la recuperación multimodal,
la interpretabilidad conceptual y el cribado de seguridad; por el
contrario, cada tarea presenta un candidato óptimo diferente.

Esta observación cuestiona la búsqueda de un único modelo dermatológico
universal y favorece un paradigma alternativo basado en la composición
de módulos especializados. En este escenario, los modelos fundacionales
actúan como infraestructura común sobre la que se construyen componentes
específicos para segmentación, clasificación, recuperación documental,
razonamiento multimodal o interpretación clínica.

La arquitectura propuesta para DermApIxel se alinea con esta visión:
más que un modelo único, constituye un ecosistema de modelos
especializados coordinados alrededor de una ontología común y de una
representación dermatológica compartida. Los resultados obtenidos
sugieren que esta aproximación modular puede ser una vía más realista y
escalable para la integración de la inteligencia artificial en la
práctica dermatológica que la búsqueda de un único sistema capaz de
resolver todas las tareas clínicas simultáneamente.

\section{Hallazgos no anticipados}
\label{sec:hallazgos}

Cinco observaciones del trabajo no se anticipaban desde la literatura y,
a diferencia de las conclusiones robustas que se sintetizan en el
capítulo~\ref{cap:limitaciones}, contradicen una expectativa previa
razonable. Por ese motivo merecen atención específica como aportaciones
empíricas con entidad propia.

\paragraph{Sensibilidad extrema al tokenizador y al \emph{prompt}.}
La AUROC \emph{zero-shot} de la familia DermLIP sobre HAM10000 varía en un
rango comparable al de cambiar de paradigma completo, sin modificar los
pesos del encoder visual, en función del tokenizador y de la formulación
de los \emph{prompts} (sección~\ref{sec:disc-zs}). Esta sensibilidad no
está cuantificada de forma sistemática en la literatura previa sobre
DermLIP y desplaza parte de la atención metodológica del modelo hacia su
configuración operativa.

\paragraph{Retrieval competitivo en la cola larga L3.}
La sustitución del clasificador lineal por \emph{retrieval} $k$-NN con
índice FAISS sobre los mismos \emph{embeddings} PanDerm Large supera al
sondeo lineal en el régimen de cola larga severa L3 de DermapixelAI~1.0
(sección~\ref{sec:dermapixel-extra-faiss}). El hallazgo indica que, cuando
el número de muestras por clase es muy bajo, un esquema no paramétrico
basado en vecindad puede ser preferible al ajuste de una frontera de
decisión paramétrica.

\paragraph{Sondeo lineal supera al ajuste fino en DDI.}
Sobre DDI ($N_{\mathrm{train}} = 427$), el sondeo lineal sobre PanDerm
Large supera al ajuste fino en accuracy, AUROC, BAcc y F1 ponderada
(sección~\ref{sec:disc-lp-vs-ft}). El comportamiento, atribuible a
sobreajuste bajo $N$ reducido, no se documenta de forma explícita en la
literatura sobre PanDerm y ofrece una heurística operativa para datasets
pequeños.

\paragraph{Más datos externos no siempre mejoran el modelo.}
La transición v4$\to$v5 de Dermapixel R0 constituye un resultado negativo
de interés: añadir $15\,128$ imágenes externas genéricas
(PAD-UFES-20, DermNet) con \emph{captions} sintéticas \emph{reduce} la
\emph{accuracy}@1 \emph{zero-shot} L3 ($0{,}2475 \to 0{,}2079$,
sección~\ref{sec:spanderm-escala-especificidad}). El resultado contradice
la hipótesis razonable de que aumentar el volumen de datos mejora el
rendimiento y sugiere que, en un dominio clínico de alta especificidad
lingüística, la alineación entre dominio y lenguaje puede ser más
determinante que la escala bruta de imágenes.

\paragraph{Interpretabilidad emergente del SAE Large sobre
Derm7pt.}
La validación independiente del SAE Large ($16\,384$
\emph{features}) sobre el dataset Derm7pt~\cite{kawahara2019},
reportada en sección~\ref{sec:sae-derm7pt}, identifica
\emph{features} individuales con AUROC efectivo hasta $0{,}823$
sobre \emph{vascular structures} y hasta $0{,}926$ sobre vasos
en horquilla ($N_{+} = 15$ positivos). Las anotaciones del
\emph{Seven-Point Checklist} no se usaron ni en el preentrenamiento
del SAE ni en el del encoder PanDerm Large. Esto implica que la base
sparse aprendida sobre Derm1M ($507\,657$ imágenes) ya ha
capturado representaciones alineadas con la taxonomía dermatoscópica
clásica. El hallazgo sostiene la interpretabilidad mediante Sparse
Autoencoders más allá de los $48$ conceptos visuales de
SkinCon~\cite{skincon} sobre los que se diseñó el experimento
original del Anexo~\ref{cap:anexof}. El \emph{Concept Bottleneck
Model} con cuello de botella sobre los siete criterios cuantifica
además el coste de esta interpretabilidad sobre la tarea binaria
melanoma/no-melanoma ($-5{,}8$\,pp AUROC, $0{,}890 \to
0{,}832$). La jerarquía de pesos del meta-clasificador
(\emph{blue whitish veil} $+2{,}54$, \emph{regression structures}
$+2{,}07$) es coherente con el peso clínico de estos criterios en
la escala diagnóstica de referencia~\cite{kawahara2019}. La lectura
operativa es que el SAE no es sólo un mecanismo de reducción
dimensional: admite interpretación concepto-a-concepto sobre el
vocabulario dermatoscópico. Esto refuerza la viabilidad del módulo de
interpretabilidad del prototipo (Anexo~\ref{cap:anexoh}) sin
reentrenar el SAE.

\section{Limitaciones interpretativas}
\label{sec:disc-limitaciones}

Conviene cerrar la discusión acotando hasta dónde llegan estas
conclusiones. Seis elementos condicionan la interpretación de los
resultados, y el capítulo~\ref{cap:limitaciones} los desarrolla en
detalle. Primero, el solapamiento exacto de Dermnet con el dataset de
preentrenamiento Derm1M ($100\%$ por hash MD5) introduce una
incertidumbre acotada sobre las cifras de la familia DermLIP en ese
dataset ---HIBA y MSKCC, pese a su origen ISIC-archive, no presentan
solapamiento MD5---; las cifras de PanDerm, DINOv2, ConvNeXt-Large,
EfficientNetV2-Large y SigLIP-Large no se ven afectadas. Conviene
además precisar el alcance de esa auditoría: el hash MD5 detecta
coincidencias exactas pero no \emph{near-duplicates} ---recortes,
reescalados o recompresiones de una misma imagen---, de modo que lo
verificado es la ausencia de solapamiento \emph{exacto}, no de toda
contaminación visual posible (sección~\ref{sec:lim-leakage}). Segundo, los
tamaños muestrales de los subgrupos Fitzpatrick V y VI ($169$ y
$73$ imágenes respectivamente en el conjunto de test) producen
intervalos de confianza amplios para las cifras de \emph{gap}, lo que
limita la cuantificación exacta de la desigualdad por fototipo.
Tercero, la sensibilidad al \emph{prompt} en zero-shot y en la
evaluación de LLMs multimodales hace que las cifras reportadas sean
condicionales a una formulación concreta; generalizarlas a otras
formulaciones requiere validación específica. Cuarto, el trabajo
evalúa rendimiento sobre fotografías de archivos de investigación; la
generalización a imagen real adquirida en flujos de teledermatología
sobre cohorte hospitalaria prospectiva no se aborda y constituye una
línea futura inmediata. Quinto, la denominación de Dermapixel R0 como
primer modelo fundacional dermatológico en español es una afirmación
de prioridad deliberadamente acotada: designa la adaptación de un
modelo fundacional preexistente al castellano, no un preentrenamiento
a gran escala, y su condición de \emph{primero} se circunscribe a esa
categoría ---pesos abiertos, texto clínico narrativo y registro en
español--- (sección~\ref{sec:linea-spanderm}); una reivindicación de
prioridad histórica exhaustiva excedería el alcance bibliográfico de
este trabajo. Sexto, la evaluación sobre DermapixelAI~1.0 se realiza
sobre un conjunto de tamaño moderado y con una distribución fuertemente
desbalanceada en los niveles ontológicos más finos: en particular, la
cola larga de L3 limita la estabilidad de las estimaciones por clase y
obliga a interpretar los resultados de ese nivel como evidencia
exploratoria más que como una ordenación definitiva entre métodos. Esta
limitación afecta especialmente a las comparaciones entre clasificadores
sobre L3 ---\emph{retrieval} $k$-NN, sondeo lineal y Dermapixel R0--- y
motiva la necesidad de ampliar el banco hospitalario en trabajos futuros
(sección~\ref{sec:lim-dermapixel}).

Estas limitaciones acotan el dominio de validez de las cifras, pero
no invalidan los patrones cualitativos identificados: la ventaja del
preentrenamiento de dominio, la saturación temprana del sondeo
lineal, el régimen $N_{\mathrm{train}} \lesssim 500$ como umbral de
preferencia del LP, la brecha entre especialistas y LLMs
multimodales generalistas, y la existencia de un \emph{gap} por
fototipo sistemático y consistente entre encoders.
